\begin{lemma}
\label{editorial-source-lemma-unique-prime-over-localize-below}
Let $R \to S$ be a ring map.
Let $\mathfrak p \subset R$ be a prime.
Assume that
\begin{enumerate}
\item there exists a unique prime $\mathfrak q \subset S$ lying over
$\mathfrak p$, and
\item either
\begin{enumerate}
\item going up holds for $R \to S$, or
\item going down holds for $R \to S$ and there is at most one prime
of $S$ above every prime of $R$.
\end{enumerate}
\end{enumerate}
Then $S_{\mathfrak p} = S_{\mathfrak q}$.
\end{lemma}

\begin{proof}
Write the original ring map as $\varphi:R\to S$.
The rings in the conclusion are the specified
localizations
$$
S_{\mathfrak p}=\varphi(R\setminus\mathfrak p)^{-1}S,
\qquad S_{\mathfrak q}=(S\setminus\mathfrak q)^{-1}S.
$$
Since $\varphi^{-1}(\mathfrak q)=\mathfrak p$,
the first multiplicative set is contained in the
second. Thus there is an original localization map
$j:S_{\mathfrak p}\to S_{\mathfrak q}$, sending
$s/\varphi(r)$ to the same fraction.
No injectivity of $\varphi$ or of these localization
maps is assumed.

A prime of $S_{\mathfrak p}$ corresponds exactly
to a prime $\mathfrak q'\subset S$ whose contraction
$\mathfrak p'=\varphi^{-1}(\mathfrak q')$ is contained
in $\mathfrak p$. A prime of $S_{\mathfrak q}$
corresponds exactly to a prime
$\mathfrak q'\subseteq\mathfrak q$. Both statements
follow from contraction and extension under
localization: the corresponding prime consists of
the fractions with numerator in $\mathfrak q'$,
and these maps are inverse because the denominator
set is disjoint from that prime.

Assume first (1) and (2)(a). For such a
$\mathfrak q'$ over $\mathfrak p'\subseteq\mathfrak p$,
going up supplies $\mathfrak q''\supseteq\mathfrak q'$
whose contraction is $\mathfrak p$. Uniqueness over
$\mathfrak p$ makes $\mathfrak q''=\mathfrak q$,
so $\mathfrak q'\subseteq\mathfrak q$.
Assume instead (1) and (2)(b). Going down applied
to $\mathfrak q$ gives
$\mathfrak q''\subseteq\mathfrak q$ over
$\mathfrak p'$. Uniqueness over $\mathfrak p'$
makes $\mathfrak q''=\mathfrak q'$, with the same
conclusion. This argument uses uniqueness only
over the primes $\mathfrak p'\subseteq\mathfrak p$;
the global uniqueness assumption in (2)(b) can
therefore be weakened to precisely those primes.

It follows in either case that for every
$u\in S\setminus\mathfrak q$, the element $u/1$
belongs to no prime of $S_{\mathfrak p}$: every
such prime contracts to a prime contained in
$\mathfrak q$. Hence $u/1$ is a unit. Indeed a
nonunit generates a proper ideal lying in a
maximal ideal, which is prime, a contradiction.
Write its inverse as $v/\varphi(r)$ with
$r\notin\mathfrak p$. The equality
$(u/1)(v/\varphi(r))=1$ means that for some
$t\notin\mathfrak p$,
$$
\varphi(t)(uv-\varphi(r))=0\quad\hbox{in }S.
$$
Thus the inverse to $j$ has the full fraction
formula
$$
k:S_{\mathfrak q}\longrightarrow S_{\mathfrak p},
\qquad
\frac su\longmapsto\frac{sv}{\varphi(r)},
$$
where $v,r$ satisfy that original inverse identity.
Independence of this choice follows from uniqueness
of the inverse of $u/1$. Independence of fraction
representatives can also be checked explicitly:
if $w(u's-us')=0$ for $w\notin\mathfrak q$,
then $w/1,u/1,u'/1$ are all units in
$S_{\mathfrak p}$, so multiplication by their
inverses gives $s/u=s'/u'$ there. Addition and
multiplication follow by the same common-denominator
identities. Consequently $k$ is the localization
extension of the original map $S\to S_{\mathfrak p}$.
Both composites fix $S$ and all the specified inverse
denominators, hence are identities. This proves
$S_{\mathfrak p}\simeq S_{\mathfrak q}$ through
the original canonical map, giving the equality
notation in the statement its exact meaning.

In fact for any original $\varphi$, $\mathfrak p$
and $\mathfrak q$ with
$\varphi^{-1}(\mathfrak q)=\mathfrak p$, the
following conditions are equivalent:
$$
\begin{gathered}
j:S_{\mathfrak p}\longrightarrow S_{\mathfrak q}
\text{ is an isomorphism};\\
u/1\text{ is a unit in }S_{\mathfrak p}
\text{ for every }u\notin\mathfrak q;\\
\varphi^{-1}(\mathfrak q')\subseteq\mathfrak p
\Longrightarrow\mathfrak q'\subseteq\mathfrak q
\text{ for every prime }\mathfrak q'\subset S.
\end{gathered}
$$
The first implies the second because these elements
are units in $S_{\mathfrak q}$. The second implies
the third since a prime of $S_{\mathfrak p}$ cannot
contain a unit; contraction then avoids every
element outside $\mathfrak q$. The third implies
the second by the maximal-ideal argument, and the
second supplies the proved inverse fraction map.
This is the exact criterion used by both cases
of the original lemma. It is specific to the
specified further localization; a bijective
spectrum map alone has not been substituted for
an isomorphism of rings.
\end{proof}